Soit $\Omega$ un cercle de centre $M$ et $\Gamma$ un cercle de centre
$N$ tels que le rayon de $\Omega$ soit strictement plus petit que le
rayon de $\Gamma$. On suppose que les cercles $\Omega$ et $\Gamma$ se
coupent en deux points $A$ et $B$ distincts. La droite $(MN)$ coupe
$\Omega$ en un point $C$ et $\Gamma$ en un point $D$ de sorte que les
points $C$, $M$, $N$ et $D$ soient alignés dans cet ordre. Soit $P$ le
centre du cercle circonscrit au triangle $ACD$. La droite $(AP)$
recoupe $\Omega$ en un point $E$ distinct de $A$ ; elle recoupe
également $\Gamma$ en un point $F$ distinct de $A$. Enfin, soit $H$
l'orthocentre du triangle $PMN$.

Démontrer que la parallèle à $(AP)$ passant par $H$ est tangente au
cercle circonscrit au triangle $BEF$.

(L'\emph{orthocentre} d'un triangle est le point d'intersection de ses
hauteurs.)
